\section{Related Works}
\paragraph{Deep Generative Models for Tabular Data Generation.}
Generative models for tabular data have become increasingly important and have widespread applications~\cite {privacy, TabCSDI, privacy_health}. To deal with the imbalanced categorical features, \citet{ctgan} proposes CTGAN and TVAE based on the popular Generative Adversarial Networks~\citep{gan} and VAE~\citep{vae}, respectively. Multiple advanced methods have been proposed for synthetic tabular data generation in the past year. Specifically, GOGGLE~\citep{goggle} became the first to explicitly model the dependency relationship between columns, proposing a VAE-based model using graph neural networks as the encoder and decoder models. Inspired by the success of large language models in modeling the distribution of natural languages, GReaT transformed each row in a table into a natural sentence and learned sentence-level distributions using auto-regressive GPT2. In recent years, the physical diffusion process has inspired a lot of advanced research in deep learning. For example, DIFFormer~\citep{difformer} develops a scalable Transformer model for geometric data via a constrained diffusion process, and the Denoising Diffusion models have achieved great success in image generation~\citep{ddpm}. STaSy~\citep{stasy}, TabDDPM~\citep{tabddpm}, and CoDi~\citep{codi} concurrently applied the popular diffusion-based generative models for synthetic tabular data generation.
\paragraph{Generative Modeling in the Latent Space.}
While generative models in the data space have achieved significant success, latent generative models have demonstrated several advantages, including more compact and disentangled representations, robustness to noise, and greater flexibility in controlling generated styles~\citep{vq-vae, vq-vae-2, vq-gan}. 
For example, the recent GAN literature~\citep{latent_gan} has demonstrated superior controllability via adversarial learning in the latent space. 
Recently, the Latent Diffusion Models (LDM)~\citep{ldm, sgm} have achieved great success in image generation as they exhibit better scaling properties and expressivity than the vanilla diffusion models in the data space~\citep{ddpm, vp, edm}. The success of LDMs in image generation has also inspired their applications in video~\citep{latent-video} and audio data~\citep{latent-audio}. To the best of our knowledge, the proposed work is the first to explore the application of latent diffusion models for general tabular data generation tasks.
